\subsection{The arithmetic quotient for polynomial cutoffs}
\label{acfd:native-consumer}

The maps in \eqref{eq:mobius-L} and \eqref{eq:spine-map} are exactly
\eqref{acfd:mobius} and $E_N=L_NP_N$.
The Hilbert space is $H_k=L^2(\Delta_k,dt;\mathbb R)$ in both constructions.
The fiber operators integrate with respect to length, so their quadratic
form is the same $\langle F,K_kF\rangle$.

\begin{corollary}[Arithmetic quotients with positive widths]
\label{acfd:native-polynomial}
Fix an admissible set of distinct nonnegative shifts $\mathcal H$.
Assume the distribution hypothesis \eqref{eq:lod} at a fixed exponent
$0<\vartheta\leq1$.
Fix $0<\delta<\vartheta/2$, put $\rho=\vartheta/2-\delta$, and use
the corresponding $W_N,R_N,C_N$ and admissible class $v_0(N)$.
Let $f$ be a fixed nonzero real polynomial on $\Delta_k$.
For all sufficiently large $N$, choose
\[
 0<\varepsilon_N<\frac1{2(k+1)},\qquad
 \varepsilon_N\longrightarrow0,\qquad
 D_0(N)\varepsilon_N^2\longrightarrow\infty,
 \qquad \lambda_N=L_NP_N(\Xi_{\varepsilon_N}f).
\]
Then $\lambda_N\in\Lambda_N$ is defined, $S_1(\lambda_N)>0$ eventually,
and
\[
 \frac{S_2(\lambda_N)}{S_1(\lambda_N)}
 =\rho R_k(f)+
 O_f\left(\sqrt{\varepsilon_N}+\frac1{D_0(N)\varepsilon_N^2}\right).
\]
The positive sequence
\[
 \varepsilon_N=\min\left\{\frac1{4(k+1)},D_0(N)^{-2/5}\right\}
\]
gives error $O_f(D_0(N)^{-1/5})$.
If $\rho R_k(f)>m$ for an integer $m\geq0$, then each sufficiently
large interval $I_N$ contains an integer with at least $m+1$ prime translates.
\end{corollary}

\begin{proof}
Every specified width is positive and lies in the domain of
Lemma~\ref{acfd:cutoff}.
That lemma proves $\Xi_{\varepsilon_N}f\in\mathcal X_k$ for each such $N$.
Hence point sampling, $L_N$, and the incidence map all have their stated domains.
Theorem~\ref{thm:forms} gives \eqref{acfd:uniform-mass} and
\eqref{acfd:uniform-prime}, with constants independent of $N$ and the
selected smooth function.  Its prime normalization is $\rho C_N$.
Theorem~\ref{acfd:moving-cutoff} therefore gives the quotient estimate
and eventual positivity.  Proposition~\ref{prop:finite} gives the prime count.
\end{proof}

For a polynomial with $R_k(f)>2m/\vartheta$, one may first choose
$0<\delta<\vartheta/2-m/R_k(f)$.
This interval is nonempty, and the choice gives $\rho R_k(f)>m$.
The polynomial and exponent loss are fixed before selecting the positive
width sequence.  For the polynomial thresholds in
Corollary~\ref{bridge:cor:prime}, the same construction applies with $m=1$.
It retains their distribution hypothesis and gives distinct qualifying
integers by taking disjoint intervals.
